\section{Scope, limitations, and conclusion}
\label{sec:conclusion}

For a complete cyclic core, the unrestricted one-level feeder converts a
single transient choice into an integer allocation of $d$ phase-rooted
subtrees.  Although the feeder state itself occupies only one vertex, the
allocation it selects is replicated across the boundary-dominant levels.
This gives the exact max--min \eqref{eq:one-level-optimizer}.  The
constant-convolution criterion then separates equality with the spectral
mean from mere improvement, and the canonical forced chain replaces the
denominator $d$ by the exact grid denominator $d^L$.

The complete-cyclic core formula is a supporting calculation.  The residual
content begins with the transient optimizer, continues with its saturation
arithmetic and canonical-chain convergence, and culminates in the Hausdorff
cyclic-essential-SCC obstruction \eqref{eq:four-state-values}.  The proofs are based on
matching cylinder and mass-distribution estimates, so the full dimension
statements do not rely on a spectral upper bound or a finite computation.

The scope is exact and limited.  The principal families are the complete
cyclic core with the unrestricted one-level feeder and the canonical
unrestricted $L$-level forced chain.  Only explicitly declared
finite-composition one-level variants are covered by
\cref{rem:finite-composition}; restricted multilevel variants require the
balanced-access hypothesis \eqref{eq:balanced-access}.  We exclude arbitrary
finite strictly transient feeder shapes, incomplete phase blocks, return
edges, nontransient reuse, and arbitrary reducible matrices.  We also do not
claim divisibility necessary without full Fourier support and do not import
version-sensitive irreducible equality clauses from the literature.

Two concrete directions remain.  First, selected non-complete cyclic blocks
may admit exact phase summaries even though label choices cease to be
equiprobable; identifying such summaries would require a new lower measure as
well as an upper count.  Second, controlled communication or return edges
would replace the finite-union decomposition by an infinite recurrence.
Any extension in that direction needs a matched upper and lower mechanism,
not only a larger finite optimizer.  Finally, the bounded source search used
for positioning does not establish priority, and no such claim is made.
